\section{Plane Isometries}
\textbf{Conventions.} Permutations compose from right to left. The identity is $e$. A group generated by an element $a$ is denoted $\gen{a}$.

\exercise{12.1}{Symmetries depend on ambient dimension}
\begin{enumerate}[label=(\alph*)]
\item For a point $c\in\R$, the isometries fixing it are $x\mapsto x$ and $x\mapsto2c-x$, the reflection through $c$. An isometry of the line has form $x\mapsto x+t$ or $x\mapsto-x+t$, and the condition of fixing $c$ determines $t$ in each case.
\item For a point $P$ in the plane, all rotations about $P$ and all reflections across lines through $P$ fix it, together with the identity. There are infinitely many, unlike in the line. Translations and genuine glide reflections fix no points.
\item A nondegenerate segment in $\R$ has only the identity and reflection through its midpoint: its two endpoints must either both stay fixed or be interchanged.
\item In $\R^2$ the symmetries are the identity, reflection in the line containing the segment, reflection in its perpendicular bisector, and the half-turn about its midpoint. They form a Klein four-group. Fixing both endpoints allows reflection in the supporting line, and interchanging them allows the other reflection or the half-turn.
\item In $\R^3$ examples include every rotation about the supporting line, reflection in any plane containing that line, and reflection in the perpendicular midpoint plane. More fully, with midpoint at the origin and segment on the $x$-axis, all maps $(x,y,z)\mapsto(\varepsilon x,Q(y,z))$ with $\varepsilon=\pm1$ and $Q$ any orthogonal $2\times2$ matrix preserve the segment.
\end{enumerate}

\exercise{12.2}{Orientation under composition}
Assign sign $+1$ to an orientation-preserving map $P$ and $-1$ to an orientation-reversing map $R$. Under composition these signs multiply, as can be seen from the determinants of the orthogonal linear parts of the isometries. Thus
\[\begin{array}{c|cc}&P&R\\\hline P&P&R\\R&R&P\end{array}\]
In particular two reversals restore orientation, while one reversal changes it.

\exercise{12.3}{Types of products of plane isometries}
Use $t$ for a nonzero translation, $p$ for a nonidentity rotation, $\mu$ for a reflection, and $\gamma$ for a genuine glide reflection. With identity outcomes omitted, every possible type is recorded by
\[\begin{array}{c|cccc}&t&p&\mu&\gamma\\\hline
 t&t&p&\mu,\gamma&\mu,\gamma\\
 p&p&t,p&\mu,\gamma&\mu,\gamma\\
 \mu&\mu,\gamma&\mu,\gamma&t,p&t,p\\
 \gamma&\mu,\gamma&\mu,\gamma&t,p&t,p
\end{array}\]
Here is why the lists are exhaustive. Orientation rules eliminate $t,p$ from a product with exactly one reversing factor, and eliminate $\mu,\gamma$ from a product with zero or two reversing factors. Translations compose by adding vectors. A translation times a rotation keeps the same nonzero rotational angle, hence is another rotation. Two rotations can add to a nonzero angle (rotation) or cancel their angles about different centers (translation). Two reflections across intersecting lines give a rotation; across distinct parallel lines they give a translation. Adding a translation to an orientation-reversing map can introduce or cancel displacement along its axis, giving either a reflection or a glide reflection. Choosing intersecting versus parallel axes for products involving reversing factors gives respectively the rotation and translation alternatives. These constructions realize every entry; inverse choices may give the omitted identity.

\exercise{12.4}{A figure with only the identity symmetry}
Take a scalene triangle, for example vertices $A=(0,0)$, $B=(4,0)$, $C=(1,2)$. Its three side lengths are $4,\sqrt5,\sqrt{13}$, all distinct. Any symmetry must preserve side lengths, so it fixes all three vertices and is the identity. A sketch is
\begin{center}\setlength{\unitlength}{5mm}\begin{picture}(5,3)
\put(0,0){\line(1,0){4}}\put(0,0){\line(1,2){1}}\put(1,2){\line(3,-2){3}}
\put(0,-0.5){$A$}\put(4,-0.5){$B$}\put(0.6,2.2){$C$}
\end{picture}\end{center}

\exercise{12.5}{A figure with two symmetries}
Use an isosceles triangle that is not equilateral, such as vertices $(-2,0),(2,0),(0,3)$. The unique vertex between the two equal sides must be fixed; the base endpoints can either stay fixed or interchange. Thus the only symmetries are the identity and reflection across the vertical altitude, forming $\Z_2$.
\begin{center}\setlength{\unitlength}{5mm}\begin{picture}(5,4)
\put(0,0){\line(1,0){4}}\put(0,0){\line(2,3){2}}\put(2,3){\line(2,-3){2}}\put(2,-0.3){\line(0,1){3.7}}
\end{picture}\end{center}

\exercise{12.6}{A figure with three symmetries}
Draw an equilateral triangle and mark a heavy dot one quarter of the way along each edge, always proceeding counterclockwise around the boundary. The dots are part of the figure.
\begin{center}\setlength{\unitlength}{5mm}\begin{picture}(5,4)
\put(0,0){\line(1,0){4}}\put(0,0){\line(1,1.732){2}}\put(2,3.464){\line(1,-1.732){2}}
\put(1,0){\circle*{0.16}}\put(3.5,0.866){\circle*{0.16}}\put(1.5,2.598){\circle*{0.16}}
\end{picture}\end{center}
The rotations through $0^\circ,120^\circ,240^\circ$ preserve the marks. A reflection reverses the boundary direction and would put the marks three quarters of the way along each counterclockwise edge, where there are no marks. Since the convex hull is the equilateral triangle, every possible symmetry was already among its rotations and reflections. Thus exactly the three rotations remain, giving $\Z_3$.

\exercise{12.7}{A figure with cyclic order-four symmetry}
Use a square with one heavy dot one quarter of the way along each edge, in a consistent counterclockwise direction:
\begin{center}\setlength{\unitlength}{5mm}\begin{picture}(5,5)
\put(0,0){\framebox(4,4){}}\put(1,0){\circle*{0.16}}\put(4,1){\circle*{0.16}}\put(3,4){\circle*{0.16}}\put(0,3){\circle*{0.16}}
\end{picture}\end{center}
All four quarter-turn rotations preserve the figure. Every square reflection reverses the direction along its boundary and moves each quarter-edge marker to a missing three-quarter-edge marker. Thus none is a symmetry. The convex hull forces every symmetry to be a square symmetry, so the group consists of the four rotations and is $\Z_4$.

\exercise{12.8}{A figure with Klein-four symmetry}
A rectangle with unequal side lengths has the identity, reflection in each of its two midlines, and the half-turn about its center. A quarter-turn would interchange its unequal dimensions, so is not allowed. Each of its three nonidentity symmetries has order $2$, and the two reflections multiply to the half-turn, giving the Klein group $V$.
\begin{center}\setlength{\unitlength}{5mm}\begin{picture}(7,4)\put(0,0){\framebox(6,3){}}\put(3,0){\line(0,1){3}}\put(0,1.5){\line(1,0){6}}\end{picture}\end{center}

\exercise{12.9}{Orders of the four isometry types}
A nonzero translation has infinite order, because its $k$th power translates by $k$ times its nonzero vector. A rotation through angle $\theta$ has finite order $q\ge2$ when $\theta/(2\pi)=p/q$ in lowest terms, and infinite order when this ratio is irrational. Every reflection has order $2$. A genuine glide reflection has infinite order: if its nonzero displacement along the axis is $v$, its square is the nonzero translation by $2v$. An element whose square has infinite order cannot have finite order.

\exercise{12.10}{Which types have fixed points?}
A nonidentity rotation fixes its center. A reflection fixes every point of its axis. A nonzero translation has no fixed point because $P+v=P$ would force $v=0$. A genuine glide reflection also has none: its component parallel to the axis moves every point by the same nonzero amount, which cannot vanish. Thus exactly rotations and reflections can have fixed points.

\exercise{12.11}{Exactly one fixed point}
Only a nonidentity rotation. Its center is fixed, and any other point moves around a circle by a nonzero angle. A reflection has an entire line of fixed points, while translations and genuine glide reflections have none.

\exercise{12.12}{Exactly two fixed points}
None of the four nonidentity types. A rotation has one, a reflection has infinitely many, and translations and glide reflections have none. The identity fixes the whole plane and also does not have exactly two.

\exercise{12.13}{Infinitely many fixed points}
Among the four nonidentity types, only reflections: their fixed-point set is their entire axis. If the identity is also allowed, it fixes every point of the plane and therefore is another isometry with infinitely many fixed points.

\exercise{12.14}{Three noncollinear fixed points force the identity}
Let $A,B,C$ be the fixed noncollinear points and let $X$ be any point. An isometry preserves its distances to all three: the image $Y$ must satisfy $YA=XA$, $YB=XB$, and $YC=XC$. The first two circles meet in at most two points, which, when distinct, are mirror images across $AB$. Since $C$ is off $AB$, those two mirror points have different distances to $C$; hence the third distance selects $X$ uniquely. Thus $Y=X$ for every $X$ and the isometry is the identity. This is the geometric reason three noncollinear distance references determine a plane point.

\exercise{12.15}{Isometries agreeing on three points}
Suppose $\phi(P_i)=\psi(P_i)$ at the three given noncollinear points. Form the isometry $\eta=\psi^{-1}\phi$. Then $\eta(P_i)=\psi^{-1}(\psi(P_i))=P_i$ for each $i$. By Exercise 14, $\eta$ is the identity. Multiplying $\psi^{-1}\phi=\mathrm{id}$ on the left by $\psi$ yields $\phi=\psi$.

\exercise{12.16}{All rotations together with the identity}
They do not form a subgroup. Let $r_A(X)=2A-X$ and $r_B(X)=2B-X$ be half-turns about two distinct centers. Then
\[(r_Br_A)(X)=2B-(2A-X)=X+2(B-A),\]
a nonzero translation. This product is neither a rotation nor the identity, so closure fails.

\exercise{12.17}{Translations together with the identity}
Yes. Write $T_v(X)=X+v$, allowing $v=0$ for the identity. Then $T_vT_w=T_{v+w}$ and $T_v^{-1}=T_{-v}$. The sums and negatives of plane vectors are plane vectors, so the set contains the identity and is closed under products and inverses.

\exercise{12.18}{Rotations about a common point}
Yes. After translating coordinates so that the fixed center is the origin, composition of rotations adds angles, and inversion negates the angle. Thus $R_\alpha R_\beta=R_{\alpha+\beta}$ and $R_\alpha^{-1}=R_{-\alpha}$, with angles read modulo $2\pi$. Angle $0$ gives the identity, so these maps form a subgroup.

\exercise{12.19}{One reflection together with the identity}
Yes. If $s$ is reflection across the specified line, then $s^2=e$. Hence $\{e,s\}$ is closed under multiplication, and each element is its own inverse. This is a subgroup isomorphic to $\Z_2$.

\exercise{12.20}{All glide reflections together with the identity}
No. A genuine glide reflection $g$ has nonzero displacement $v$ along its axis, and $g^2$ is translation by $2v$. This is neither a genuine glide reflection nor the identity. Thus the product of an allowed element with itself leaves the proposed set.

\exercise{12.21}{Isometries in finite subgroups}
Every element of a finite group has finite order. Nonzero translations and genuine glide reflections have infinite order and cannot occur. Reflections can occur with order $2$, and rotations can occur when their angle is a rational multiple of $2\pi$. Thus the possible nonidentity types are precisely reflections and finite-order rotations.

\exercise{12.22}{The rotation subgroup of a finite plane group}
Restrict the orientation sign map to $G$: its value is $+1$ on preserving elements and $-1$ on reversing elements, and signs multiply. Its $+1$ set $H$ is therefore closed under products and inverses and contains the identity. In a finite group, the preceding exercise excludes nonzero translations, so $H$ consists exactly of the rotations and identity described in the question. If there is no reversing element, $H=G$. Otherwise fix a reversing $s\in G$; multiplication by $s$ is a bijection from $H$ to the reversing elements, because it changes the sign and has inverse multiplication by $s^{-1}$. Hence $|G|=2|H|$.

\exercise{12.23}{A finite group of rotations about a point is cyclic}
If $G=\{e\}$, it is already cyclic. Otherwise choose the smallest positive angle $\theta\in(0,2\pi)$ occurring among its finitely many rotations, and let $r$ be that rotation. For any rotation angle $\alpha\in[0,2\pi)$, division by $\theta$ writes $\alpha=q\theta+\delta$ with $0\le\delta<\theta$. The rotation $r^{-q}R_\alpha$ lies in $G$ and has angle $\delta$. Minimality of $\theta$ forces $\delta=0$, so every element is a power of $r$. This proves $G=\gen r$. Since its order is finite, $r$ also has finite order and its multiples close up exactly around the circle.

\exercise{12.24}{Repeated F figures}
\textbf{Convention.} Here ``nontrivial glide'' means one not expressible as a product of a translation and reflection that are themselves symmetries, as defined on printed page 116.
\begin{enumerate}[label=(\alph*)]\item No.\item No.\item No.\item No.\item The group is isomorphic to $\Z$.\end{enumerate}The upper and lower parts of F differ, and its left and right sides differ. Only integral translations by the spacing preserve the repeated upright figures.

\exercise{12.25}{Repeated T figures}
\textbf{Convention.} Here ``nontrivial glide'' means one not expressible as a product of a translation and reflection that are themselves symmetries, as defined on printed page 116.
\begin{enumerate}[label=(\alph*)]\item No.\item No.\item Yes.\item No.\item The group is isomorphic to $D_\infty$.\end{enumerate}Each T has a vertical mirror, but no horizontal mirror or half-turn symmetry. A unit translation $t$ and a vertical reflection $s$ generate the group, with $s^2=e$ and $sts=t^{-1}$.

\exercise{12.26}{Repeated E figures}
\textbf{Convention.} Here ``nontrivial glide'' means one not expressible as a product of a translation and reflection that are themselves symmetries, as defined on printed page 116.
\begin{enumerate}[label=(\alph*)]\item No.\item Yes.\item No.\item No.\item The group is isomorphic to $\Z\times\Z_2$.\end{enumerate}Horizontal reflection through the middle bar preserves every E; vertical reflection and half-turn do not. This reflection commutes with the unit translation. Glides obtained by combining these two symmetries are trivial in the book's sense.

\exercise{12.27}{Repeated Z figures}
\textbf{Convention.} Here ``nontrivial glide'' means one not expressible as a product of a translation and reflection that are themselves symmetries, as defined on printed page 116.
\begin{enumerate}[label=(\alph*)]\item Yes.\item No.\item No.\item No.\item The group is isomorphic to $D_\infty$.\end{enumerate}A half-turn about the center of a Z preserves it, but neither horizontal nor vertical reflection does. The half-turn $r$ reverses translation: $rtr=t^{-1}$, yielding the infinite dihedral group.

\exercise{12.28}{Repeated H figures}
\textbf{Convention.} Here ``nontrivial glide'' means one not expressible as a product of a translation and reflection that are themselves symmetries, as defined on printed page 116.
\begin{enumerate}[label=(\alph*)]\item Yes.\item Yes.\item Yes.\item No.\item The group is isomorphic to $D_\infty\times\Z_2$.\end{enumerate}The H has both horizontal and vertical mirrors and a half-turn. Translation and a vertical mirror generate $D_\infty$, while the horizontal mirror is an independent commuting involution. Every glide is the product of an existing translation and a horizontal reflection.

\exercise{12.29}{Alternating bent strokes}
\textbf{Convention.} Here ``nontrivial glide'' means one not expressible as a product of a translation and reflection that are themselves symmetries, as defined on printed page 116.
\begin{enumerate}[label=(\alph*)]\item No.\item No.\item No.\item Yes.\item The group is isomorphic to $\Z$.\end{enumerate}The upper and lower bent strokes are horizontal mirror images displaced by half a repeat. A glide $g$ through that half spacing interchanges the two rows, and $g^2$ is the full translation. No pure reflection or half-turn matches the direction of the short horizontal stubs. All symmetries are powers of $g$.

\exercise{12.30}{Alternating cup and cap figures}
\textbf{Convention.} Here ``nontrivial glide'' means one not expressible as a product of a translation and reflection that are themselves symmetries, as defined on printed page 116.
\begin{enumerate}[label=(\alph*)]\item Yes.\item No.\item Yes.\item Yes.\item The group is isomorphic to $D_\infty$.\end{enumerate}Vertical lines through cup or cap centers are mirrors. A half-period glide $g$ in the horizontal baseline interchanges cups and caps; its reflection component alone is not a symmetry. Together with a vertical reflection $s$, it gives $sgs=g^{-1}$, so the group is infinite dihedral. Their product also yields half-turns.

\exercise{12.31}{The square grid}
\textbf{(a)} Rotations through $90^\circ$ and $180^\circ$ occur, about square centers or appropriate grid vertices; no smaller positive angle is possible because the two perpendicular edge directions must be preserved as a set. \textbf{(b)} Yes: horizontal, vertical, and diagonal mirror lines preserve the grid. \textbf{(c)} No nontrivial glides in the book's sense. A symmetry has the form $X\mapsto AX+v$ with $v$ a lattice vector and $A\in D_4$ a symmetry fixing a grid vertex. If it is a glide, $A$ is a reflection, so it is already a product of a translation and a reflection in the symmetry group. Ordinary glides do exist; the word ``nontrivial'' is essential here.

\exercise{12.32}{The staggered square or brick pattern}
\textbf{(a)} A half-turn is a symmetry, but a quarter-turn is not: vertical grid lines continue indefinitely, whereas horizontal edges terminate at the staggered joins. Thus the only angle in $(0,180^\circ]$ is $180^\circ$. \textbf{(b)} Yes. Take unit square tiles in column $k$ shifted vertically by $k/2$. There are horizontal mirrors at integer or half-integer heights and vertical mirrors through column centers. \textbf{(c)} No nontrivial glides under the book's product definition. The group is generated by its centered rectangular translation lattice and the two perpendicular reflections through a column center on a horizontal mirror. Every glide has one of those reflection linear parts and differs from such a reflection by a translation already in the group. The translation/reflection factors need not be the canonical parallel-axis factors; this is why specifying the book's definition matters.

\exercise{12.33}{Square tiles marked L}
\textbf{(a)} No nonidentity rotation preserves the identically upright L markings. \textbf{(b)} No mirror preserves these markings. \textbf{(c)} No nontrivial glide occurs: an orientation reversal would turn each L into a reflected letter, which no translation restores. Thus only the square-lattice translations remain. As usual for a letter-decorated pattern, the intended L is the ordinary nonsymmetric letterform in the source, not a specially designed equal-arm right-angle symbol with an extra diagonal mirror.

\exercise{12.34}{Square tiles marked E}
\textbf{(a)} No nonidentity rotation preserves upright E. \textbf{(b)} Yes: reflection in a horizontal line through a row of letter centers (and suitable midway lines) preserves the pattern. Vertical or diagonal reflections do not. \textbf{(c)} No nontrivial glide in the source's sense: every glide has horizontal reflection linear part, and composing it with a horizontal mirror symmetry yields a translation symmetry. Hence it is a product of existing translation and reflection symmetries.

\exercise{12.35}{Square tiles marked H}
\textbf{(a)} The angle $180^\circ$ occurs, but $90^\circ$ turns each H sideways and is not allowed. \textbf{(b)} Yes: horizontal and vertical mirror lines through suitable rows and columns preserve H. \textbf{(c)} No nontrivial glide: both possible reflection directions already occur as mirrors. Composing any glide with a parallel mirror symmetry gives a translation symmetry, expressing the glide as a product of symmetries of these two types.

\exercise{12.36}{The regular hexagonal tiling}
\textbf{(a)} The possible angles in $(0,180^\circ]$ are $60^\circ,120^\circ,180^\circ$. Rotations about hexagon centers by multiples of $60^\circ$ preserve the tiling, and a symmetry must permute its three unoriented edge directions, excluding other angles. \textbf{(b)} Yes: lines through opposite vertices or opposite side midpoints of a hexagon are mirrors of the entire tiling. \textbf{(c)} No nontrivial glide in the book's sense. Taking a hexagon center as origin, each symmetry is a lattice translation followed by an element of the hexagon's dihedral point group; when reversing orientation, the latter is an existing reflection. Thus any glide already has the required factorization.

\exercise{12.37}{Hexagons marked with upright triangles}
\textbf{(a)} The angle $120^\circ$ occurs; $60^\circ$ and $180^\circ$ turn the centered upward triangle to a downward one, so they fail. \textbf{(b)} Yes: the three reflection axes of each upward equilateral triangle also preserve the hexagonal tiling. \textbf{(c)} No nontrivial glides. The remaining point group is $D_3$, and every orientation-reversing linear part is one of those existing mirrors; composing a glide with that mirror leaves a translation in the same symmetry group. This gives precisely the product excluded by the word ``nontrivial.''

\exercise{12.38}{Horsemen tessellation}
Ignore color shading as the instructions require, but preserve the outlines and the internal horse-and-rider markings.
\textbf{(a)} Half-turns ($180^\circ$) occur; quarter-turns and smaller rotations do not, since a horizontal horse would be turned into an absent vertical orientation. \textbf{(b)} There are no reflection axes: a mirror reverses the handed rider/horse outline and its offset adjoining shapes. \textbf{(c)} There are nontrivial glides in both horizontal and vertical directions. Reflecting and shifting by half the appropriate period matches adjacent oppositely oriented rows or columns; the reflection alone does not preserve the pattern. \textbf{(d)} Translation generators may be chosen horizontal and vertical, with directions $(1,0)$ and $(0,1)$. Choose the lengths as the smallest full repeats of a marked upright horse. The half-repeat displacements attached to the glides are not pure translations. These observations describe the usual $pgg$ pattern.

\exercise{12.39}{Imaginary human figures tessellation}
\textbf{(a)} Rotations through $90^\circ$ and $180^\circ$ occur. Quarter-turn centers are the junctions where four corresponding feet form a pinwheel; a quarter-turn carries one figure orientation to the next. \textbf{(b)} Yes: horizontal and vertical lines along the bilateral body axes are mirrors. There are no diagonal mirrors. Candidate axes must preserve the repeating junction pattern. In one full repeat, diagonal candidates through a four-foot junction reverse its chiral pinwheel, while shifted diagonal candidates between junction rows fail to match the neighboring head and body orientations. Horizontal and vertical mirrors can exchange different junctions, so this exclusion does not assert that every pinwheel has the same handedness. \textbf{(c)} Yes, there are nontrivial diagonal glide reflections. To see why, compose a quarter-turn symmetry with a horizontal mirror symmetry. The result reverses orientation and has the linear reflection part for a diagonal axis. Since no diagonal reflection is a symmetry, this composition must be a genuine glide reflection. It cannot be a product of a translation symmetry and a reflection symmetry: a translation does not change the linear part, so such a factorization would require an absent diagonal mirror. This distinguishes the pattern ($p4g$) from an ordinary square grid. \textbf{(d)} Translation generators have diagonal directions $(-1,1)$ and $(1,1)$, with lengths taken to the first repeat of a figure in the same orientation. Following these diagonal chains, rather than moving to a differently rotated neighboring figure, identifies the translation periods. Color differences are disregarded. The printed odd-answer entry for part (c) says ``No''; the diagonal glides just constructed give ``Yes'' under the section's own definition.

\exercise{12.40}{Reptile tessellation}
\textbf{(a)} Ignoring shading, rotations through $60^\circ,120^\circ,180^\circ$ occur. The six reptile orientations cycle around the sixfold junctions; rotations through $120^\circ$ also occur at the threefold junctions. \textbf{(b)} No: a reflection reverses the handed arrangement of the curled bodies, feet, and tails, and that reversed arrangement is absent. \textbf{(c)} No glide reflections occur, since a glide also reverses orientation and cannot restore the handed reptile markings. \textbf{(d)} Choose two primitive translation directions along rows of identically oriented reptiles in the triangular lattice, for example $(1,0)$ and $(1,\sqrt3)$, with suitable lengths to reach the next full repeat. These directions have angle $60^\circ$; the directions, rather than a scale read from the printed image, are requested. The uncolored pattern has point group $C_6$ (wallpaper type $p6$).

\exercise{12.41}{Human figures tessellation}
\textbf{(a)} The possible nonzero angle at most $180^\circ$ is $120^\circ$, rotating the three figure orientations around their threefold junctions. A $60^\circ$ turn would require an orientation not present. \textbf{(b)} Yes: the bilateral human figures align along mirror axes in the pattern. \textbf{(c)} No nontrivial glides in the book's sense; an existing reflection with the same linear part removes the reversing component and leaves a translation symmetry. \textbf{(d)} Two translation-generator directions are $(0,1)$ and $(1,\sqrt3)$; choose their lengths to carry a figure to the next identically oriented one. These are directions only, as requested, so the lengths need not be inferred from the reproduction. Shading distinctions are not retained.

\exercise{12.42}{Cube rotations and $S_4$}
\idea Let a cube rotation act on the four long (unoriented) diagonals through its center.
This produces a homomorphism $\Phi$ into $S_4$, because performing two rotations in succession composes their diagonal permutations. We show its kernel is trivial. If a rotation preserves all four diagonal lines, choose unit vectors $v_1,\ldots,v_4$ along them directed to the vertices of one inscribed regular tetrahedron; then $v_i\cdot v_j=-1/3$ for $i\ne j$ and $v_1+\cdots+v_4=0$. The rotation sends $v_i$ to $\varepsilon_i v_i$. Preservation of the nonzero pairwise dot products gives $\varepsilon_i\varepsilon_j=1$ for every pair, so all signs agree. If all are $+1$, three independent diagonals are fixed and the rotation is the identity. If all are $-1$, the linear map is $-I_3$, which has determinant $-1$ and is not a rotation. Thus $\Phi$ is injective.
There are $24$ cube rotations: choose a bottom face in six ways and a front face in four ways. Since $|S_4|=24$, the injective map is also surjective and gives an isomorphism. It is important to use lines rather than just one endpoint of each diagonal, and to rule out central inversion as above.
